% !TeX root = ../thesis.tex
\chapter{文字层级与版面间距}
\label{chap:introduction}

\section{模板的验证目标}
学位论文模板不仅要给出命令，还应让使用者直观看到命令的最终效果。本章把标题、正文、段落和行距集中在连续页面中，以便检查字体是否正确嵌入、标题层级是否清晰、段前段后距离是否稳定，以及同一字号在不同软件中的视觉密度是否接近。

\section{标题层级}
章标题采用三号黑体居中；一级节标题采用四号黑体；二级节标题采用小四号黑体；更低层级分别使用小四号楷体。标题前后的距离由类文件统一控制，不建议在正文中用空行或手工的 \verb|\vspace| 调整。

\subsection{二级标题示例}
二级标题用于组织同一节内相互关联的论述。标题后第一段与后续段落采用一致的首行缩进，避免依赖编辑器自动插入空白字符。

\subsubsection{三级标题示例}
三级标题适合标识局部方法或分析步骤。若标题层级过深，通常意味着章节结构可以重新组织。

\subsubsubsection{四级标题示例}
四级标题为行内标题，适合简短条目。正式论文应尽量保持各章层级深度一致。

\section{正文行高与 Word 多倍行距}
字号、行高和行距是三个不同概念。字号描述字形的标称尺寸；行高或基线间距描述相邻两行基线之间的距离；行距则是字形上下边界之间剩余的空白。Word 的“单倍行距”通常已经大于一倍字号，“多倍行距 1.3”是在该内部基准上继续放大，而不是简单计算为字号的 1.3 倍。LaTeX 的 \verb|\linespread| 同样会乘在字体预设的基线距离上，因此不能把两个软件中的倍率直接视为同一个物理量。关于这些概念和换算关系的进一步说明见文献~\cite{wdlobin2026typesetting}。

本模板没有继续叠加全局 \verb|\linespread|，而是在 \verb|\normalsize| 中把小四号正文设置为 12 bp 字号和 20 bp 基线间距，即 \verb|\@setfontsize\normalsize{12bp}{20bp}|。比值约为 1.67。这个结果用于逼近学校规范中 Word 多倍行距 1.3 的视觉密度；不同 Word 版本、字体度量和打印驱动仍可能产生细微差异，因此应使用同一段文字进行实页比较。

\subsection{整页连续文本检验}
下面的连续文字用于观察一页 A4 纸范围内的基线稳定性。排版比较时，应关闭 PDF 阅读器的“适合宽度”缩放，分别以 100\% 比例查看 PDF 与 Word 导出的 PDF。比较对象应使用同一纸张大小、页边距、字体、字号和段落缩进。不要只比较屏幕截图中两行文字的距离，因为截图缩放和抗锯齿会改变观感。

科技论文的正文通常由定义、推理、证据和结论构成。清晰的段落应围绕一个中心问题展开：开头说明本段讨论的对象，中间给出必要的数据或推导，末尾说明这些信息如何支持本节论点。排版系统负责建立一致的视觉节奏，但不能替代内容组织。稳定的行高使读者在长行之间移动视线时不易串行，适当的段落长度则帮助读者识别论证边界。

比较行距时还要关注中英文混排。中文字符大多接近方形字面，英文小写字母具有升部和降部，公式中的积分号、上下标及分式还可能超过正文常规字面范围。固定基线距离必须为这些元素保留足够空间。如果行距过小，多行公式、带上下标的变量或加粗标题会显得拥挤；如果行距过大，段落内部的联系会减弱，页面也会出现不必要的疏松感。

段前段后距应与行高分开判断。本模板正文段落不额外增加段间距，而用首行缩进表示新段落；标题则通过类文件中的 \verb|\titlespacing| 和 \verb|\ctexset| 设置上下距离。图表与正文之间的距离由 \verb|\textfloatsep|、\verb|\intextsep| 和 \verb|\floatsep| 控制。多行公式内部的额外距离由 \verb|\jot| 控制。这些参数职责不同，混用会导致短页面看似正常、长页面却出现不一致的空白。

实际校准建议以学校提供的 Word 规范页作为参照。先在 Word 中使用指定字体、小四号和多倍行距 1.3，导出 PDF；再把同一段纯文本放入本节，分别测量二十行或三十行文字从第一条基线到最后一条基线的总高度。若累计差异明显，只调整类文件中 \verb|\normalsize| 的第二个参数，例如在 19.5 bp 至 20.5 bp 范围内试验，而不要同时修改字号和段落距离。

当正文中出现公式、图表或列表时，页面节奏会受到局部环境影响。因此，行距测试既需要纯文本页，也需要包含典型元素的综合页面。下一章给出符号和公式，第三章给出表格、算法、代码与图形，可共同用于检查模板在真实科技论文场景中的稳定性。
